\subsubsection{Recursos disponibles}

El desarrollo del prototipo se llevará a cabo con recursos propios del
equipo de trabajo, conformado por tres integrantes. A continuación se
describen los recursos disponibles para el proyecto.

\textbf{Recursos humanos.} El equipo está integrado por tres
desarrolladores estudiantes de Ingeniería en Sistemas Computacionales,
quienes aportarán su tiempo y conocimiento en las áreas de desarrollo
móvil, backend, infraestructura en la nube y aprendizaje automático.
El trabajo se realizará de forma colaborativa utilizando equipos de
cómputo personales de cada integrante.

\textbf{Recursos de hardware.} Los recursos físicos disponibles para
el proyecto son los siguientes:

\begin{itemize}
    \item \textbf{Equipos de cómputo personales:} Cada integrante
    cuenta con su propio equipo para el desarrollo, entrenamiento del
    modelo y pruebas del sistema.
    \item \textbf{Servidor local Raspberry Pi con disco duro externo:}
    Adquirido con un costo de \$5,000 MXN, se utilizará como entorno
    de pruebas local para validar el comportamiento del backend antes
    del despliegue en la nube.
    \item \textbf{Dron:} Adquirido con un costo de \$3,000 MXN, se
    utilizará para la captura aérea de imágenes del arbolado del
    Bosque de Chapultepec durante las jornadas de recolección del
    conjunto de datos.
    \item \textbf{Dispositivo Android:} Disponible para pruebas de
    la aplicación móvil en condiciones reales de campo.
\end{itemize}

\textbf{Recursos de software.} Todas las herramientas de desarrollo
seleccionadas son de código abierto o cuentan con licencias gratuitas
para uso académico, incluyendo Android Studio, IntelliJ IDEA Community,
Python, TensorFlow, OpenCV, Scikit-learn, Angular CLI y GitHub Actions.
Esto elimina cualquier costo asociado a licencias de software durante
el desarrollo del prototipo.
